%%%%%%%%%%%%%%%%%%%%%%%%%%%%%%%%%%%%%%%%%
% Medium Length Professional CV
% LaTeX Template
% Version 2.0 (8/5/13)
%
% This template has been downloaded from:
% http://www.LaTeXTemplates.com
%
% Original author:
% Rishi Shah 
%
% Important note:
% This template requires the resume.cls file to be in the same directory as the
% .tex file. The resume.cls file provides the resume style used for structuring the
% document.
%
%%%%%%%%%%%%%%%%%%%%%%%%%%%%%%%%%%%%%%%%%

%----------------------------------------------------------------------------------------
%	PACKAGES AND OTHER DOCUMENT CONFIGURATIONS
%----------------------------------------------------------------------------------------
\renewcommand{\baselinestretch}{0.99}
\documentclass{resume} % Use the custom resume.cls style

\newenvironment{indentpar}[1]%
  {\begin{list}{}%
          {\setlength{\leftmargin}{#1}}%
          \item[]%
  }
  {\end{list}}

\usepackage[T1]{fontenc}
\usepackage[utf8]{inputenc}
\usepackage{mathptmx}
\usepackage[]{natbib}
\usepackage[left=1in,top=1in,right=1in,bottom=1in]{geometry} % Document margins
\usepackage{datetime2}

\linespread{1}
\fontsize{12}{13.5}
\newcommand{\tab}[1]{\hspace{.25\textwidth}\rlap{#1}}
\newcommand{\itab}[1]{\hspace{0em}\rlap{#1}}
\renewcommand{\bibsection}{}

% Date command
\usepackage[absolute]{textpos}
\usepackage[UKenglish]{isodate}
\setlength{\TPHorizModule}{1mm}
\setlength{\TPVertModule}{1mm}
\newcommand{\lastupdated}{\begin{textblock}{60}(155,5)
\fontsize{8pt}{10pt}\selectfont 
Last Updated \today
\end{textblock}}

%%===================================================================================================
%%===================================================================================================

\name{YILI YANG} % Your name
\address{+44 7490334233\textbar\textbar yyang@woodwellclimate.org\textbar\textbar ORCID 0000-0002-1791-3899} % Your phone number and email

%-------------------------------------------------------------
%====================================================
\begin{document}
% Derived metrics generated from data/stats.json + cv/publications.bib by
% scripts/emit-metrics-tex.js. Do not hardcode these numbers — edit the source data.
% cv/**/generated/metrics.tex
% GENERATED by scripts/emit-metrics-tex.js — DO NOT EDIT BY HAND.
% Source: data/stats.json + cv/publications.bib. Regenerate with `npm run metrics`.
% Last generated: 2026-08-12T12:29:34.345Z
\newcommand{\journalpapers}{13}
\newcommand{\conferences}{8}
\newcommand{\datasets}{1}
\newcommand{\pubcount}{20}
\newcommand{\citations}{291}
\newcommand{\hindex}{8}
\newcommand{\itenindex}{7}
\newcommand{\commits}{693}
\newcommand{\repos}{22}
\newcommand{\invitedtalks}{17}
\newcommand{\peerreviews}{8}
\newcommand{\nsfreviews}{1}
\newcommand{\mentees}{9}
\newcommand{\funding}{99,749}

\lastupdated
\begin{rSection}{Research SCOPE}
My research develops/applies advanced AI methodologies to solve complex scientific problems, especially in geoscience, environmental and climate sciences. I create data/AI-driven approaches that accelerate scientific discovery, improve predictions, and transform our understanding of our changing planet through interdisciplinary collaborations.

\end{rSection}
%====================================================
\begin{rSection}{Education}
{\bf University of Edinburgh, UK} 

\begin{indentpar}{0.5cm}
\textbf{PhD in Petrophysics} \hfill {Oct 2016 - Apr 2021}

\textbf{MSc (Research) in Geology, Distinction} \hfill {August 2015 - August 2016}

\textbf{BSc Geology} \hfill {Sep 2011 - May 2015} 
\end{indentpar}

\end{rSection}
%====================================================
\begin{rSection}{Experience}
% Experience entries are single-sourced from data/experience.json and rendered by
% scripts/render-cv.js (templates/biosketch-experience.tex.njk). Edit the data, not here.
% GENERATED by scripts/render-cv.js — DO NOT EDIT BY HAND.
% Source: data/experience.json via cv/templates/biosketch-experience.tex.njk. Run `npm run render`.

{\bf AI Scientist -- Geospatial \& Remote Sensing, Woodwell Climate Research Center, MA, USA (remote)} \hfill {Jan 2022 - }

\begin{itemize}
    \item Project Lead: the pan-Arctic Retrogressive Thaw Slumps (RTS) mapping initiative, applying state-of-the-art deep learning techniques to satellite imagery for detecting and quantifying permafrost thaw features, directly contributing to climate feedback modelling and carbon cycle research
    \item Designed and built the ARTS scientific dataset, establishing a benchmark training resource for deep learning applications in environmental monitoring
    \item Published the first circumpolar retrogressive thaw slump probability map and polygon inventory (NSF Arctic Data Center, doi:10.18739/A26970107), produced by a 41.5-million-tile deep learning inference sweep of 2025 PlanetScope imagery: 19,068 quality-controlled slump polygons covering 529.7 km$^2$, released as cloud-optimised rasters and open vectors
    \item Developed innovative machine learning workflows for multimodal data integration, combining various satellite imagery types and field observations to enhance feature detection accuracy
    \item Developed A Zero-shot, interactive segmentation workflow for satellite image labelling using the Segment Anything Model (SAM)
    \item Collaborated with UCONN in using Location Embeddings to improve satellite image segmentation with vision foundation models.
    \item Collaborated with ASU in using big geospatial data to produce pan-Arctic statistics.
    \item Collaborated with UIUC in optimising the workflow of the Deep Learning model inference in big geospatial data (Arctic-scale) in GCP.
    \item Collaborated with ASU in modelling abrupt thaw initialisation with remote sensing and deep neural networks.
    \item Collaborated in two wildfire projects using deep learning for wildfire mapping on satellite imagery and using environmental variables to predict wildfire using machine learning
    \item Collaborated in carbon flux time series imputation using machine learning
    \item Led the machine learning workshop series for Woodwell researchers in 2022 and 2023, providing training in advanced AI techniques for environmental science applications including image processing and time-series forecasting. Designed progressive curriculum materials and adapted teaching methods for diverse learner backgrounds.
    \item Mentored interns and research assistants, created training materials for interns and research assistants. Mentee Works were presented at the AGU annual conference.
\end{itemize}

{\bf Data Science Fellow, Faculty.ai, London, UK} \hfill {Sep 2021 - Dec 2021}

\begin{itemize}
    \item Completed intensive fellowship in Machine Learning and Artificial Intelligence, focusing on developing practical AI solutions for complex data challenges. Include trainings and courses for professional data scientist and an industry placement.
\end{itemize}

{\bf PhD Researcher, International Centre for Carbonate Reservoirs, Edinburgh, UK} \hfill {Sep 2017 - Sep 2020}

\begin{itemize}
    \item Funded by PetroBras, petrophysics and digital rock experimental study.
    \item Conducted doctoral research on multiphase fluid flow in porous media, developing novel computational methods for analyzing complex flow dynamics, funded by Petróleo Brasileiro
    \item Parcipitated in petrophysical experiments at the Diamond Light Source UK, the Advanced Photon Source Argonne National Laboratory, US and the Swiss Light Source Paul Scherrer Institute, CH
\end{itemize}


\end{rSection}
%===============================================================
\begin{rSection}{PUBLICATIONS, CONFERENCES AND DATA SETS}
\nocite{yang2026keeping}
\nocite{virkkala2026data}
\nocite{perera2025panarctic}
\nocite{potter2026burned}
\nocite{clelland2025machine}
\nocite{li2025multiscale}
\nocite{yang2025collaborative}
\nocite{yang2024arctic}
\nocite{gu4762408multi}
\nocite{li2024segment}
\nocite{rodenhizer2024comparison}
\nocite{yang2024arts}
\nocite{rogers2023permafrost}
\nocite{yang2023mappingb}
\nocite{potter2023mapping}
\nocite{potter4803815mapping}
\nocite{rodenhizer2023comparison}
\nocite{li2023geoai}
\nocite{yang2023mappinga}
\nocite{mullen2023using}
\nocite{singh2022mapping}
\nocite{yang2022fast}
\nocite{marti2021time}
\nocite{marti2020chemical}
\nocite{yang2018immiscible}
\bibliographystyle{unsrtnat}
\bibliography{../publications}
\end{rSection}
%=====================================================
\begin{rSection}{RESEARCH FUNDINGS}

Funding for Climate Solutions: A Generic Climate AI Framework for Multi-domain Time Series Prediction, Woodwell Climate Research Center, \$\funding
    
Overall Objectives: This project will attempt to develop a generic climate AI framework that will unify and accelerate time-series inferences across Woodwell Climate Research Center's diverse research domains, from Arctic carbon fluxes to tropical forest dynamics.

Pending/ongoing/submitted funding applications with NASA, NSF, Google.org and the Bezos Foundation
    
\end{rSection}
%=====================================================
\begin{rSection}{PARTICIPATED RESEARCH INITIATIVES}

\textbf{Permafrost Discovery Gateways} Funded by Google.org and NSF, the Gateway is making information of permafrost conditions available throughout the Arctic by providing access to big geospatial products and tools to allow exploration and discovery for researchers, educators, and the public at large. This includes providing automated monthly monitoring of permafrost thaw during the snow-free season from satellite imagery.

\textbf{Permafrost Pathways} Funded through the TED Audacious Project — a collaborative funding initiative catalyzing big, bold solutions to the world’s most urgent challenges. Through a joint effort between Woodwell Climate Research Center, the Arctic Initiative at Harvard Kennedy School, and the Alaska Institute for Justice, Permafrost Pathways brings together leading experts in climate science, policy action, and environmental justice to inform and develop adaptation and mitigation strategies to address permafrost thaw.

\textbf{GeoAI Challenge} Cooperate with the Cyberinfrastructure and Computational Intelligence Lab and Prof. Wenwen Li at Arizona State University to host an NSF-funded AI challenge (Cyber2A  Arctic Challenge) using remote sensing imagery to detect permafrost thaw.

\end{rSection}
%=====================================================
\begin{rSection}{Mentorship}
Multi-domain Timeseries Forecasting 
\begin{indentpar}{0.5cm}
Sandeep Poudel \hfill2026\\
Dogukan Teber \hfill2025
\end{indentpar}

Satellite Image Labelling Workflow and Training Image Generation
\begin{indentpar}{0.5cm}
Robb Young \hfill2025\\
Marly Bartkus \hfill2025
\end{indentpar}

Deep Learning label data curation for ESRI-Woodwell Collaboration
\begin{indentpar}{0.5cm}
Marly Bartkus \hfill2025
\end{indentpar}

Automatic Improvement of Geospatial Digitisation using Segment Anything Model
\begin{indentpar}{0.5cm}
Arav Agarwal \hfill2025\\
Milagros Becerra \hfill2025
\end{indentpar}

Estimation of the Volumetric Loss of Retrogressive Thaw Slumps using DEM
\begin{indentpar}{0.5cm}
Margo Moceyunas \hfill2024
\end{indentpar}

Using Vision-Transformer to segment Retrogressive Thaw Slumps on Remote Sensing Imagery
\begin{indentpar}{0.5cm}
Ridhima Singh  \hfill2023
\end{indentpar}
\end{rSection}
%=====================================================
\begin{rSection}{ACADEMIC SERVICES, CONFERENCES AND TALKS}
\begin{itemize}
    \item Manuscripts review: \peerreviews
    \item NSF Proposal review: \nsfreviews
    \item Tropical Forest Forever Facility (TFFF): Pathways to Impact, London Climate Action Week, UK \hfill 2026
    \item Accelerating Climate Progress with AI: From Science to Action Workshop, The National Academies, US \hfill 2024
    \item Invited talk at the Arctic Carbon Flux Upscaling Workshop, Sweden \hfill 2024
    \item Invited panel talk at the Esri UC Spatial Analytics Summit, US \hfill 2024
    \item Invited talk at the Cyber2A Workshop, US \hfill 2024
    \item European Geophysical Union Annual Conference, EU \hfill 2022, 2023, 2024
    \item Google Geo-for-Good Summit, US/Singapore \hfill 2022, 2023, 2025
    \item Tri-Polar Remote Sensing Conference, China \hfill 2023, 2024
    \item RTSInTrain Workshop, EU \hfill 2023
    \item NASA Arctic-Boreal Vulnerability Experiment workshop, US \hfill 2022
    \item Invited Permafrost Discovery Gateway Webinar \hfill 2022
    \item Invited webinar at the Royal Meteorological Society: Machine Learning for Atmospheric Sciences: Values and Controversies, UK \hfill 2022
    \item American Geophysical Union annual conference, US \hfill 2018
\end{itemize}


\end{rSection}
%=====================================================
\end{document}
